\section{Rationale}

Cognitively engaged physical activity and exergaming may influence executive outcomes in children with ADHD, but the available evidence does not imply that adding movement to every cognitive task improves transfer \parencite{benzing2019exergaming,zhao2024cognitivephysical,fang2024cognitivephysical,pradana2026exergaming}. Movement changes arousal, motor demand, visual input, response modality, and often task novelty at the same time.

Augmented reality introduces further factors: camera use, spatial anchoring, real-environment search, device motion, physical navigation, and environmental variability. Pediatric AR studies therefore make feasibility, tolerability, and safety central outcomes rather than assuming that AR is intrinsically advantageous \parencite{welsby2025ar}.

The ARC-ADHD companion study asks whether contextual variation can support generalization of a stable regulation rule. The target rule remains constant:

\[
\text{challenge}
\rightarrow
\text{pause}
\rightarrow
\text{strategy}
\rightarrow
\text{recovery}
\rightarrow
\text{return to goal}.
\]

What changes is the sensorimotor and environmental context in which the rule is executed.

The central claim is deliberately weaker than the claim that AR improves ADHD. The study tests whether embodiment and AR alter proximal control and whether repeated practice across contexts predicts better transfer to independent tasks.
